\section{Background}
\subsection{Pathophysiological Links Between AD and T2DM}
AD and T2DM share several disease mechanisms that play out at both the molecular and whole-body at scales. Chronic insulin resistance, the defining feature of T2DM, contributes to AD by interfering with how neurons take up glucose and respond to insulin signals in the brain \cite{kciuk2024alzheimer}. This neurodegenerative change: the combination of excess glucose and disrupted insulin signaling fuels inflammation and oxidative stress, which then accelerates the formation of $\beta$-amyloid plaques and tau tangles (the signature brain abnormalities we see in AD) \cite{kciuk2024alzheimer}.

As established by the KG, Apolipoprotein $\epsilon$4 (APOE $\epsilon$4) offers a genetic link between these two diseases, example of a triple that shows this link is "Individuals with T2DM have APOE4-related cognitive and olfactory impairment". While we know APOE $\epsilon$4 as the strongest genetic risk factor for late-onset AD, its influence extends beyond just amyloid processing; it's also tied to how our bodies handle fats and glucose \cite{yassine2020apoe}. People with diabetes who carry APOE $\epsilon$4 experience faster mental decline and face higher dementia risk than diabetics without this gene variant, suggesting that the allele and metabolic stress amplifies each other's effects \cite{kciuk2024alzheimer, jabeen2022genetic}.

Inflammation represents another critical connection point between T2DM and AD. In T2DM, visceral fat and insulin-resistant tissues pump out elevated levels of adipokines and inflammatory cytokines. These molecules can breach the blood-brain barrier, adding fuel to the neuroinflammation already present in AD. From a clinical standpoint, T2DM patients carry roughly twice the risk of developing vascular dementia and show higher rates of Alzheimer's dementia across many studies \cite{kciuk2024alzheimer, exalto2012update}.

\subsection{Knowledge Graphs in Biomedical Research}
The complex web of relationships in biomedicine has pushed researchers to develop knowledge graphs (KGs) as a way to organize factual information in a structured, queryable format. Unlike traditional text documents, KGs represent biomedical entities like genes, diseases, drugs, proteins as nodes or entities in a network, with their relationships depicted as labelled edges. For example, an edge might connect a gene node to a disease node with the label ``associated with," making the relationship explicit and machine-readable. Large-scale biomedical KGs like SPOKE show how they created KGs by integrating over 40 curated databases, including DrugBank and GWAS catalogs, resulting in a comprehensive network of approximately 42 million nodes across 28 entity types and 160 million relationships \cite{soman2024biomedical}. CoDe-KG \cite{anuyah2025automated} is another SOTA pipeline built with open-source LLMs. What sets KGs apart from text-based resources is their built-in traceability. While bioinformaticians have traditionally used KGs for drug repurposing and gene discovery, we leverage them as knowledge sources for question answering (QA).

Given this context, we propose to leverage CoDe-KG \cite{anuyah2025automated}, an automated KG construction pipeline, to build focused biomedical knowledge graphs and evaluate their impact on trustworthy domain-specific questions. CoDe-KG is an open-source framework that extracts structured facts from text by combining robust co-reference resolution with syntactic decomposition. This approach breaks down complex statements and resolves pronouns, thereby capturing more complete and context-rich relations. According to the authors, CoDe-KG achieved an increase of $\Delta = $ 8\% F$_1$  when compared to prior methods on the REBEL relation extraction dataset \cite{anuyah2025automated}, and integrating co-reference + decomposition increased recall on rare relations by over 20\% \cite{ anuyah2025automated}. We harness this pipeline to distil knowledge from biomedical abstracts into three KGs of varying scope: $\mathbb{G}_1, \mathbb{G}_2$ and $\mathbb{G}_3$ for the combined AD+T2DM domain. We then use these graphs in an LLM-driven RAG setup to answer questions. By comparing performance across $\mathbb{G}_1, \mathbb{G}_2$, and $\mathbb{G}_3$, we examine whether a targeted disease-specific KG yields better answers than a broader cross-domain KG, or vice versa. We further probe how the structure and scope of knowledge graphs affect the LLM’s ability to deliver correct, well-supported answers. While RAG models often outperform non-RAG approaches, this is not always guaranteed. As \cite{gao2025frag, linders2025knowledge} point out, earlier KG-RAG frameworks with fixed search parameters could retrieve redundant trivial facts or miss important multi-hop connections, ultimately weakening the LLM's reasoning capabilities and making RAG-enhanced systems perform worse than their non-RAG counterparts.
\subsection{Our Approach}
In this work, we built three different knowledge graphs ($\mathbb{G}_1, \mathbb{G}_2$ and $\mathbb{G}_3$) to investigate how the scope of knowledge affects QA performance in this domain. We detail the curation of Probe1 and Probe2 in the subsequent sections. 
In Probe1, formulated on $\mathbb{G}_3$ and Probe2, formulated on $\mathbb{G}_1 \cap \mathbb{G}_2$, we tested different combinations of LLMs built on the different RAG systems.
Our hypothesis going in was that the focused KG ($\mathbb{G}_1$ and $\mathbb{G}_2$) might perform better on questions specifically about the AD-T2DM intersection, since it provides denser, more concentrated context for that particular overlap. Meanwhile, the merged KG ($\mathbb{G}_3$) should theoretically handle a wider range of questions about either disease or their connections more effectively. However, there's a potential downside to the merged approach as it could introduce distractors (facts relevant to one disease but not the question at hand), which might confuse the LLM. By testing the same set of questions against both knowledge graphs, we can observe any trade-offs between having a highly focused knowledge source versus a more comprehensive one.
% $\mathbb{G}_1 \cap \mathbb{G}_2$ represents the intersection of T2DM and AD knowledge, it captures only those triples with cosine similarity $\geq$ 0.65, focusing on deeply interconnected domain-specific relationships without the clutter of tangentially related topics. $\mathbb{G}_3$ takes a different approach by merging knowledge from both domains comprehensively, which allows broader cross-domain relationships to emerge (like diabetes-related genes that also appear in AD research, regardless of similarity scores).




% \section{Background}

% \subsection{KG for Finding Disease Commonalities}
% Besides being useful for finding useful patterns from raw data, KG has the potential for being especially helpful for spotting commonalities across diseases that don't seem related at first glance. Take the Alzheimer's Knowledge Base, for instance. Using a Random Forest classifier trained on topological graph features, researchers identified 8 genes shared between Alzheimer's and Parkinson's \cite{romano2024alzheimer}. These overlaps point to common neurological pathways in neurodegeneration, though it's worth noting that shared genes don't always translate to identical disease mechanisms. PrimeKG takes this further by integrating data across 10 biological scales, which allows for finding disease similarities based on shared protein perturbations, metabolic pathways, and how symptoms actually present \cite{chandak2023building}. This multi-scale perspective indicates that diseases traditionally considered distinct may, in fact, share therapeutic vulnerabilities.

% Beyond genetics, there appears to be value in simply looking at phenotypes through the lense of KG. When analyzing overlapping symptoms and clinical presentations, disease relationships start to emerge that one might miss otherwise. Systems combining the Human Phenotype Ontology with disease knowledge graphs can cluster diseases based on how similar they look clinically, and these phenotypic similarities often (though not always) reflect shared underlying mechanisms \cite{nunes2023multi, kafkas2019ontology}. Though it is not a perfect correlation, but the pattern is there.

% \subsection{KG for Disease Pattern Discovery in AD and Diabetes}
% KG is good at creating useful patterns from raw data, like connections between molecular factors and what we actually see clinically \cite{chandak2023building}. Existing studies like \cite{romano2024alzheimer, pu2023graph, yang2025alzheimer, li2024dalk} focus on KG creation for AD, whereas \cite{sousa2025gene, qin2025building, wang2020construction, cheng2025enhancing} use KG for analysing different factors related to diabetics.

% In addition, the connection between Alzheimer's and diabetes has drawn considerable interest in recent years. Knowledge graph approaches have started to reveal what look like shared biological mechanisms: insulin resistance, oxidative stress, inflammation, and protein aggregation all seem to play a role in both conditions \cite{yang2025alzheimer, almeida2022semantic, karki2017comorbidity, hu2020shared}. That said, it's not entirely clear whether these are truly common pathways or just similar-looking processes that happen to show up in both diseases. Therefore, in this paper, we will explore AD and diabetics related literature using KG, and try to look into probable commonalities between them.